\documentclass{article}
\usepackage{amsmath,amssymb,xurl}
\usepackage[hidelinks]{hyperref}
% Shared commands for notes in docs/paper-gaps/.

\DeclareUnicodeCharacter{2080}{\ifmmode{}_{0}\else\textsubscript{0}\fi}
\DeclareUnicodeCharacter{2081}{\ifmmode{}_{1}\else\textsubscript{1}\fi}
\DeclareUnicodeCharacter{2082}{\ifmmode{}_{2}\else\textsubscript{2}\fi}
\DeclareUnicodeCharacter{2099}{\ifmmode{}_{n}\else\textsubscript{n}\fi}

\newcommand{\Lean}{\textsc{Lean}}
\ExplSyntaxOn
\NewDocumentCommand{\leanid}{m}{
  \ifmmode
    \textsf{\detokenize{#1}}
  \else
    \group_begin:
    \urlstyle{sf}
    \tl_set:Nn \l_tmpa_tl {#1}
    \tl_replace_all:Nnn \l_tmpa_tl {\_} {_}
    \exp_args:NV \path \l_tmpa_tl
    \group_end:
  \fi
}
\ExplSyntaxOff
\providecommand{\tr}{\operatorname{tr}}
\newcommand{\tnleanIssueBase}{https://github.com/LionSR/TNLean/issues/}
\newcommand{\tnleanPrBase}{https://github.com/LionSR/TNLean/pull/}
\newcommand{\tnleanDocBase}{https://sirui-lu.com/TNLean/docs/}
\newcommand{\ghissue}[1]{\href{\tnleanIssueBase#1}{GitHub issue \##1}}
\newcommand{\ghpr}[1]{\href{\tnleanPrBase#1}{PR \##1}}
\newcommand{\doclink}[1]{\href{\tnleanDocBase#1.html}{\path{#1}}}

% Dirac notation and the identity operator, as in blueprint/src/macros/common.tex.
% \providecommand keeps note-local definitions authoritative.
\usepackage{dsfont}
\providecommand{\ket}[1]{|#1\rangle}
\providecommand{\bra}[1]{\langle #1|}
\providecommand{\braket}[2]{\langle #1 | #2 \rangle}
\providecommand{\ketbra}[2]{|#1\rangle\!\langle #2|}
\providecommand{\Id}{\mathds{1}}
\providecommand{\one}{\mathds{1}}
\providecommand{\C}{\mathbb{C}}
\providecommand{\MN}[1]{M_{#1}(\C)}
\providecommand{\MD}{M_D(\C)}

% Verdict marker: \gapnote{<kind>}{<status>}. Kind is the policy
% classification (clarification, local-correction, scope-restriction,
% unfaithful, false-source, open-gap); status is open, wip (elimination
% underway), resolved, or historical. Machine-read by the site index;
% prints a verdict line.
\newcommand{\gapnote}[2]{\par\noindent\textbf{Verdict:} #1 (#2).\par}

\title{Eliminating One Native Face Generator}
\author{}
\date{}
\begin{document}
\maketitle
\gapnote{open-gap}{open}

Schuch, Cirac, and P\'erez-Garc\'ia, \cite[Lemma~6.14]{Schuch2010PEPS},
source \path{Papers/1001.3807/paper_v3.tex}, lines 2199--2214,
use virtual invariance to deform a fixed-endpoint flux string.
A finite geometric classification for prescribed allowed square supports
requires further work. The auxiliary result here proves the algebraic
validity and locality of one elementary native face substitution.

Let $X=\mathbb Z_w\times\mathbb Z_h$ with $w,h>0$. In the free group on two
copies of $X$, retain one positive east generator and one positive north
generator at each native base. For the face with lower-left vertex $v$,
write $A,B,C,D$ for its bottom, right, left and top positive generators.
The reverse-ordered native face word and the two substitutions are
\begin{align}
    r_v=BAC^{-1}D^{-1},\qquad
    \sigma_v^{\rm W}(C)=D^{-1}BA,\qquad
    \sigma_v^{\rm S}(A)=B^{-1}DC.
    \label{eq:scp10_face_elimination_substitutions}
\end{align}
Each substitution fixes every generator except the displayed one.
\footnote{\leanid{TNLean.PEPS.torusDualWestFaceElimination};
\leanid{TNLean.PEPS.torusDualSouthFaceElimination}.}

If $w>1$, the left and right generators are distinct, so
$\sigma_v^{\rm W}(r_v)=1$. If $h>1$, the top and bottom generators are distinct,
so $\sigma_v^{\rm S}(r_v)=1$. No lower bound on the other positive period
is needed. At any positive periods, west substitution fixes the relation
at $u$ when $u\ne v$ and $u+(1,0)\ne v$; south substitution fixes it when
$u\ne v$ and $u+(0,1)\ne v$. These are native generator-incidence conditions.
\footnote{\leanid{TNLean.PEPS.torusDualWestFaceElimination_face};
\leanid{TNLean.PEPS.torusDualSouthFaceElimination_face};
\leanid{TNLean.PEPS.torusDualWestFaceElimination_face_eq};
\leanid{TNLean.PEPS.torusDualSouthFaceElimination_face_eq}.}

For an arbitrary group $H$ and a homomorphism $\varphi$ from this free group,
flatness on the original face implies
\begin{align}
    \varphi(r_v)=1
    \ \Longrightarrow\
    \varphi(B)\varphi(A)=\varphi(D)\varphi(C)
    \ \Longrightarrow\
    \varphi\circ\sigma_v^{\rm W}=\varphi
    \quad\hbox{and}\quad
    \varphi\circ\sigma_v^{\rm S}=\varphi.
    \label{eq:scp10_face_elimination_representations}
\end{align}
The last equalities follow by solving for the selected generator.
They need neither commutativity nor finiteness of $H$, and remain valid even
when the corresponding period equals one.
\footnote{\leanid{TNLean.PEPS.torusDualWestFaceElimination_preserves_flatRepresentation};
\leanid{TNLean.PEPS.torusDualSouthFaceElimination_preserves_flatRepresentation}.}

This establishes one local substitution, not a sequence eliminating every
allowed face. Constructing a finite descending face order, composing the
substitutions without changing previously eliminated relations, and proving
the resulting reduced-word classification for arbitrary proper support
remain separate obligations. No such classification, physical density
criterion, or endpoint-only equivalence is asserted here.
The remaining source interpretation stays under \ghissue{8703};
the Hamiltonian work under \ghissue{8947} is unaffected.

\bibliographystyle{alpha}
\bibliography{references}
\end{document}
